\documentclass[11pt]{article}

\usepackage[margin=1in]{geometry}
\usepackage{amsmath, amssymb}
\usepackage{parskip}
\usepackage[hidelinks]{hyperref}
\pagestyle{plain}

\newenvironment{reviewerquote}
  {\begin{quote}\itshape}
  {\end{quote}}

\title{Response to reviewers\\
  \large The Cooperation Ceiling: Extrinsic Population Dynamics and the Intrinsic Escape}
\author{Harry Foster, Vince Knight, Sebastian Krapohl}
\date{\today}

\begin{document}
\maketitle

Dear Dr Sasaki,

We would like to thank both reviewers for their constructive comments on our
manuscript, all of which we have addressed. Below, we have responded to each
comment, with the reviewer's comment in italics and our response following in
upright type.

\section*{Reviewer 1}

\begin{reviewerquote}
  The connection between the general theorem in the first part of the manuscript
  and the public-goods analysis in the second part could be made more explicit.
  Let N denote the population size and n the number of participants in a game.
  The theorem in the first part appears sufficiently general to apply not only
  when N = n but also when N > n. By contrast, the public-goods analysis in the
  second part is restricted to the case N = n, in which the population size and
  the number of game participants coincide. This restriction should be stated
  explicitly before the manuscript moves from the general theory to the specific
  example. 
\end{reviewerquote}

We have added a new section (section 6) which discusses the impact of
formulation of the public goods game with either \(N=n\) or \(N>n\). This
section gives a similar parameter sweep to those in sections 3.3.1 and 4.3,
applied to a public goods game where \(N>n\). We clarify in section 3.3.1 that
the initial analysis is restricted to \(N=n\), and then give results in section
6 where \(N>n\)

\begin{reviewerquote}
The manuscript is organised throughout around the contrast between extrinsic and
intrinsic dynamics, but I do not think that this single axis is sufficient. The
finding that the cooperation level does not exceed 1/2 even when r > N is
influenced not only by the use of extrinsic dynamics but also by the assumption
N = n in the public-goods game. By contrast, when small groups are repeatedly
formed from a sufficiently large population, so that N > n, the analysis is
based on expected payoffs, and cooperation can evolve when r > n. I therefore
suggest organising the discussion along two axes:
\begin{enumerate}
\item whether N = n or N > n
\item whether the update compares the payoffs of individuals who are currently present in the same state or
\end{enumerate}
compares the focal individual's current payoff with the payoff it would obtain
after changing strategy. Using these two axes would make both the position of
the manuscript and the meaning of its results clearer. 

\end{reviewerquote}

The new section 6 discusses the game along the first recommended axis, and
focuses on the impact of game formulation on the ceiling theorem. We show how
formulating a game with \(N > n\) does not yield pointwise dominance for
defectors due to the formulation of the game, and thus the ceiling can be
breached. We discuss the conditions under which cooperators are favoured, and
give an interpretation of this result as whether players consider the expected
success of their strategy across all possible scenarios, or whether they
consider the success of strategies in a realisation of the game.

\section*{Reviewer 2}

\begin{reviewerquote}
    Theorem 4 assumes that in every population state, every defector's payoff
    strictly exceeds every cooperator's payoff ("pointwise dominance of
    defection"). This is a strong condition for public goods games. Indeed, this
    condition may not be satisfied in some situations. The authors should
    clarify how this pointwise dominance condition actually holds in
    heterogeneous contribution public goods games.
\end{reviewerquote}

%TODO: address

\begin{reviewerquote}
    Prior works on aspiration/introspection dynamics (e.g., Ref. [8]) have
    already observed that intrinsic update rules can outperform imitation rules
    for promoting cooperation. I think that the authors can explicitly highlight
    the novelty of their work and the differences between their paper and these
    prior works. They may stress the general classification plus the rigorous
    upper bound theorem for extrinsic dynamics, not merely restating that
    intrinsic dynamics perform better.
\end{reviewerquote}

%TODO: address

\begin{reviewerquote}
    There are some pervious works on intrinsic dynamics, e.g., Chaos, Solitons
    and Fractals 208 (2026) 118289, IEEE Transactions on Neural Networks and
    Learning Systems 36 (2025) 8387-8400, which are closely related with this
    study, I think. These paper have demonstrated that the cooperation ceiling
    can be broken. The authors may mention these works and further compare their
    work with these papers.
\end{reviewerquote}
%TODO: address


\section*{Additional Edits}

In addition to the comments of the reviewers, we have made the following
revisions:
\end{document}